\scanpage{98}

\[
\sum|x_i x_i'|\leq N(x)N'(x')\unsure{)}\tag*{(10)}
\]

Indeed, we have $\sum_{1}^{n}|x_i x_i'|=\langle |x|_n,|x'|_n\rangle\leq N(|x|_n)N(|x'|_n)\leq N(x)N'(x')$\unsure{)}, whence (10). There results a \uline{pair-}\newline\uline{ing} between $\ell^N$ and $\ell^{N'}$, \deleted{\unsure{which} \inserted{\unsure{is}}} \inserted{whence by restricting} \uline{between $\ell_N$ and $\ell^{N'}$}; \deleted{\unsure{this pairing} \unsure{we have thus associated to $\ell$ of $\ell^{N'}$} \unsure{i.e. a map|a form} continuous linear \inserted{\unsure{of norm} $\leq 1$}, $\ell^{N'}$ into $(\ell_N)'$} \unsure{I say} that \uline{the dual of $\ell_N$ is identified exactly with $\ell^{N'}$} (with its norm) by this pairing.

\margin{\textcolor{red}{!}\quad\tikz[baseline=-.6ex]{\draw[red,line width=.7pt] (0,0)--(0,.6);}}

Indeed, \deleted{\unsure{let the}} it suffices to show that a continuous linear form $X'$ on $\ell_N$ is defined by an \inserted{$x'$} element of $\ell^{N'}$ \deleted{of \illegible{1}} such that $N'(x')\leq\lVert X'\rVert$.

Indeed, $X'$ \deleted{\unsure{is entirely} \inserted{\unsure{dé|also}}} it suffices to see it for the \deleted{\unsure{part}} restriction to $\ell^1$, then $X'$ \inserted{is} defined \deleted{\unsure{by|by}} \inserted{by} an $x'\in\ell^\infty$, \deleted{\unsure{with}} we have $N'(|x'|_n)\leq\lVert X'\rVert$, whence at once $N'(x')\leq\lVert X'\rVert$, and the result.

-- Note that in general, the dual of $\ell^N$ is not $\ell^{N'}$, \deleted{\unsure{because}} because $\ell_N$ will not be dense in $\ell^N$. Example\,: $N=N_\infty$. Note that

\margin{\uline{Proposition}}
\uline{for $\ell_N$ to be reflexive, it is necessary and sufficient that $\ell^N$ respectively $\ell_{N'}$ respectively $\ell^{N'}$ be so}. \inserted{\uline{It is necessary and sufficient for this that we have $\ell_N=\ell^N$} \uuline{and} \uline{$\ell_{N'}=\ell^{N'}$.}} \deleted{Indeed,}

$\ell^N$ reflexive implies $\ell_N$ reflexive (closed subspace) which implies $\ell^{N'}$ reflexive (dual of $\ell_N$) which \inserted{likewise} implies $\ell_{N'}$ reflexive and itself implies $\ell^N$ reflexive, which completes the circular proof \deleted{of the equivalence for} of the equivalence of the 4 conditions. Then $\ell_{N'}$, which is a \deleted{\unsure{strong}} \uline{total} \inserted{and strongly closed} subspace of the dual $\ell^{N'}$ of $\ell_N$, is identical to $\ell^{N'}$, and for the same reason we shall have $\ell_N=\ell^N$.

\margin{\textcolor{red}{?}}

\inserted{Conversely,} if this is so, then the dual of $\ell_N$ is $\ell^{N'}=\ell_{N'}$, and the dual of the latter is $\ell^N=\ell_N$, thus $\ell_N$ coincides with its bidual, and is therefore reflexive.

\inserted{Let $N$ be a Schatten norm on $R^{(N)}$ \deleted{not equivalent to $N_\infty$}\inserted{not equivalent to $N_\infty$}.} Let $E,F$ be two Hilbert spaces, \deleted{\unsure{$N$|if}} we denote by $L^{\deleted{\unsure{N}}N}_{\deleted{\unsure{0}}}(E,F)$ the \deleted{\unsure{sub|the space}} subspace of $L_0(E,F)$ formed of the $u$ such that $\unsure{(\rho_i(u))|(\rho_i u)}\in\ell^N$, i.e. such that

\[
N(u)=N(\unsure{(\rho_i(u))|(\rho_i u)})<+\infty\tag*{(11)}
\]

\margin{99}
